% !PRINTED 290
\weberpage{290}

If one composes two of these substitutions, one must observe that
\[
(\psi, \psi_b) = (\psi_a, \psi_{ab})
\]
holds, so that one has
\[
(\psi, \psi_a)(\psi, \psi_b) = (\psi, \psi_{ab})
\]
These substitutions therefore compose in entirely the same way as the subsidiary groups do according to \S.\,4 of this volume, and from this it follows:

\noindent 1. \gesperrt{The Galois group of the partial resolvent belonging to $Q$
is isomorphic with the group $P/Q$ complementary to $Q$.}

Let us now take any composition series of $P$ with the associated
index series (\S.\,6):
\weq{1}{
\begin{gathered}
P,\ P_1,\ P_2,\ \ldots,\ P_{\mu-1},\ 1 \\
j_1,\ j_2,\ \ldots,\ j_{\mu-1},\ j_\mu,
\end{gathered}
}
then, by what has just been said, the group of the equation $f = 0$ is reduced from $P$ to $P_1$
by the adjunction of \gesperrt{one} root of a normal equation of degree $j_1$. Then
it is reduced to $P_2$ by a root of a normal equation of degree $j_2$,
and so on, and finally the equation is completely solved by a root
of a normal equation of degree $j_\mu$.

On this procedure of solution a new light now falls from the theorem on the
invariability of the index series (\S.\,6, I.).

\noindent 2. \gesperrt{To solve the equation $f = 0$, one has to adjoin, one after
another, one root each of a normal equation of the degrees $j_1, j_2, \ldots, j_\mu$.
The degrees of these resolvents can indeed be altered in their order,
but not in their totality.}

The numbers $j_1, j_2, \ldots, j_\mu$ are thus far more deeply connected with the nature of an
equation, or more generally with the fields defined by the equation,
than, say, the degree of the equation; for while the degree can be
changed in manifold ways by transformation, when one introduces functions of the roots as new
unknowns, the numbers $j_1, j_2, \ldots, j_\mu$ always remain preserved and
are to be regarded as true \gesperrt{invariants} of the normal field defined by
the equation.
